% Template for post or presentation abstracts submitted to ILAS 2022
% 1. Fill in the presenter's information (name, affiliation, email
% address) and talk information (title of presentation, abstract).
% 2. Compile to make sure there are no errors.
% 3. Save the .tex file for your abstract with a distinctive name,
% ideally "FamilynameGivenname.tex".
% 4. Upload your .tex file to https://clr.ie/129557
%       
% * Please keep your LaTeX commands simple, and avoid the use of
%    user-defined macros.
% * Include details of co-authors and funding acknowledgements after the abstract.
% * Upload only the LaTeX file (with .tex extension).
% * Questions: Contact Niall Madden (Niall.Madden@NUIGalway.ie)

\documentclass[a4paper]{amsart} 
\usepackage{amssymb} 
\usepackage{hyperref}
\pagestyle{empty}
\date{} 

%% IGNORE THE NEXT 4 LINES
\makeatletter % 
\def\labelenumi{\theenumi.}
\def\theenumi{\@arabic\c@enumi}
\makeatother
%% STOP IGNORING NOW

\begin{document}

\title{On the Laplacian spectra of token graphs}
\author{Cristina Dalf\'o} %Presenting author only
\address{Universitat de Lleida} % Affiliation of presenting author
\email{cristina.dalfo@udl.cat} % Email address of presenting author only

\maketitle

% The abstract  ***no more than one page***

The $k$-token graph $F_k(G)$ of a graph $G$ is the
graph whose vertices are the $k$-subsets of vertices from $G$, two of which being adjacent whenever their symmetric difference is a pair of adjacent vertices in $G$. In this talk, we give a relationship between the Laplacian  spectra of any two token graphs of a given graph.
In particular, we show that, for any integers $h$ and $k$ such that $1\le h\le k\le \frac{n}{2}$, the Laplacian spectrum of $F_h(G)$ is contained in the Laplacian spectrum of $F_k(G)$. Besides, we obtain a relationship between the spectra of the $k$-token graph of $G$ and the $k$-token graph of its complement $\overline{G}$. This generalizes a well-known property for Laplacian eigenvalues of graphs to token graphs.

%% Coauthor details here (optional - delete if not applicable)
% and funding acknowledgment (optional - delete if not applicable)
\emph{This is joint work with Frank Duque (University of Antioquia), Ruy Fabila-Monroy (Cinvestav), Prof. Miquel Àngel Fiol (UPC), Clemens Huemer (UPC), Ana Laura Trujillo-Negrete (Cinvestav), Francisco J. Zaragoza Mart\'inez (UAM).}

\begin{thebibliography}{1}
\bibitem{my_key_for_OT49}
		K. Audenaert, C. Godsil, G. Royle, T. Rudolph.
\newblock Symmetric squares of graphs.
\newblock \emph{J. Combin. Theory B} 97:74--90, 2007.

\bibitem{cflr17}
W. Carballosa, R. Fabila-Monroy, J. Lea\~nos, L. M. Rivera. 
\newblock Regularity and planarity of token graphs.
\newblock \emph{Discuss. Math. Graph Theory} 37(3):573--586, 2017.

\bibitem{ffhhuw12}
R. Fabila-Monroy, D. Flores-Pe\~{n}aloza, C. Huemer, F. Hurtado, J. Urrutia, D. R. Wood.
\newblock Token graphs.
\newblock \emph{Graphs Combin.} 28(3):365--380, 2012.
\end{thebibliography}


\end{document}


% Please leave the next bit alone  
%%% Local Variables: 
%%% mode: latex
%%% TeX-master: "../ILAS-2022"
%%% End:
